\begin{subsection}{Ideals and Filters in Lattices} 
The subject of this section is the analog of ideals and their properties, such as being prime or principal in rings for distributive lattices. 
There are actually two structures in lattices that fit the concept: ideals and filters, and both will be relevant.
\begin{definition}
        An ideal of a semilattice \( (L, \vee) \) is a non-empty set \( I \subseteq L \) s.t.
        \begin{itemize}
            \item \( a, b \in I \implies a \vee b \in I \) 
            \item \( a \in I \) and \( b \leq a \implies b \in I \).
        \end{itemize}
        If \( L \) is bounded with element \( e \), we additionally require \( e \in I \).
        The dual notion of an ideal of a bounded semilattice is a filter \(F \subseteq L \) with \( F \neq \emptyset \), namely
        \begin{itemize}
            \item \( a, b \in F \implies a \wedge b \in F \) 
            \item \( a \in F \) and \( b \geq a \) imply \( b \in F \).
        \end{itemize}
        If \( L \) is bounded, we also require \( e \in F \).
        An ultrafilter \( U \) is a filter, that is proper and is not contained in any other proper filter.
        An ideal \( P \subset L \) is a prime ideal if its complement is a filter.
        For a bounded lattice, a set \( F \subset L \) is a filter if it is one for \( (L, \vee, \bot) \),
        similarly an ideal of a bounded lattice is one for \( (L, \wedge, \top) \).
        In particular, the second defining property for ideals is equivalent 
        to the more familiar \( a \in I, b \in L \implies a \wedge b \in I \).
    \end{definition}
    \begin{remark}
        For further use, denote by \( \text{Spec } L \) the set of prime ideals of a semilattice, also called the spectrum of \( L \).
    \end{remark}

    \begin{proposition} \label{propositionprime}
        For \( L \) a distributive lattice an ideal \( P \subset L \) being prime is equivalent to:
        \begin{itemize}
            \item  \( P \) is proper (\( P \subsetneq L \)) and \( a \wedge b \in P \implies a \in P \text{ or } b \in P \).
        \end{itemize}
        If \( L \) is also bounded, \( P \) being prime is equivalent to the following: 
        \begin{itemize} 
            \item \( \top \notin P \) and \( a \wedge b \in P \implies a \in P \text{ or } b \in P\).
            \item \( P \) is the kernel of a lattice homomorphism \( f : A \to \{0, 1\} \), where \( \{0, 1\} \) is the binary lattice.
        \end{itemize}
    \end{proposition}
    \begin{proof}
        Theorem 9.8 of \cite{birkhoff1948lattice}.
    \end{proof}
     Especially the first property reminds us of classical prime ideals from rings. 
        Standard examples of an ideals are sets of the form \( (-\infty, a] \) in \( \mathbb{R} \), and the reader may convince themself that it is also prime (yielding a specimen of a filter too). 
        Filters are a very flexible tool that can be used to give an alternative axiomatic foundation for topology by leveraging the fact that the power set forms a lattice with intersections as the meet and union as the join.
        Ultrafilters are the dual concept of maximal ideals and encode in the topology setting convergence to a point.

    \begin{proposition} \label{meetirreducibleprime}
        Let \( L \) be a distributive lattice, then the principal ideal \( (a)^\downarrow  \)
        is prime if and only if \( a \) is meet irreducible,
        meaning there are no \( b, c \in L \setminus \{a\} \) s.t. \( a = b \wedge c \).
    \end{proposition}
    \begin{proof}
        Exercise 9.6.4(a) in \cite{birkhoff1948lattice}.
    \end{proof}
    We will later study "meet irreducible" elements of \( \extrm \), although the reader is invited to already make their own guesses.
    There will be a natural interpretation, what makes them irreducible in terms of risk.

    \begin{proposition} \label{separationlatspec}
        In a distributive lattice \( L \), every ideal is contained in a prime ideal.
        Moreover prime ideals can distinguish points of the lattice, meaning for distinct \( a, b \in L \),
        there is a prime ideal \( P \) s.t. only one of \( a,b \) is in \( P \).
    \end{proposition}
    \begin{proof}
        Corollary 2.1.17 in \cite{genlattheorygraetzer}.
    \end{proof}
    \begin{remark}
        A consequence of Proposition \ref{separationlatspec} is that if we endow 
        \( \spec{L} \) with the direct analog of the Zariski topology,
        meaning a subbasis of open sets is given by \( D(a) := \{ \mathfrak{p} \in \spec{L} | a \notin \mathfrak{p} \} \)
        for \( a \in L \), referred to as Stone topology, then \( \spec{L} \) is a \( T_0 \)-space.
        In classical lattice theory, this is the basis for a series of celebrated representation theorems
        such as that of Stone and Birkhoff. We won't leverage this topology on \( \spec \extrm \), because
        there are far more prime ideals than our use cases would permit.
    \end{remark}
\end{subsection}
